\QHeading{第1题　口算。}
\par\noindent


\begin{oralgrid}
\OralItem{1}{400\times30}{12000}
\OralItem{2}{50\times200}{10000}
\OralItem{3}{130\times30}{3900}
\OralItem{4}{11\times80}{880}
\OralItem{5}{19\times40}{760}
\OralItem{6}{360\div20}{18}
\OralItem{7}{630\div90}{7}
\OralItem{8}{32\times30}{960}
\OralItem{9}{20\times15}{300}
\OralItem{10}{450\div15}{30}
\OralItem{11}{660\div11}{60}
\OralItem{12}{500\div25}{20}\end{oralgrid}
\par

\QHeading{口算思路补充（对应上表题号）}

\Answer{上面列出得数，下面按小题给出不同计算思路。计数单位法适用于小数，整数题可比较凑整、拆分或商不变等方法。}
\MethodTitle{第2小题：$50\times 200$}

\MethodBlock{口算思路一：拆因数，凑整计算}{把一个因数拆成两个因数，先凑成整百或整千。}{\begin{aligned}[t]&\phantom{={}}50\times 200\\
&=50\times (2\times 100)\\
&=(50\times 2)\times 100\\
&=100\times 100\\
&=10000\end{aligned}}

\MethodBlock{口算思路二：把特殊因数改成整百或整千}{先把25、125等改写成整百或整千除以一个数。}{\begin{aligned}[t]&\phantom{={}}50\times 200\\
&=(100\div 2)\times 200\\
&=100\times (200\div 2)\\
&=100\times 100\\
&=10000\end{aligned}}
\MethodTitle{第4小题：$11\times 80$}

\MethodBlock{口算思路一：凑整后补偿}{把接近整十、整百的因数凑整，再补回或减去多算的部分。}{\begin{aligned}[t]&\phantom{={}}11\times 80\\
&=80\times (10+1)\\
&=80\times 10+80\times 1\\
&=800+80\\
&=880\end{aligned}}
\MethodTitle{第5小题：$19\times 40$}

\MethodBlock{口算思路一：凑整后补偿}{把接近整十、整百的因数凑整，再补回或减去多算的部分。}{\begin{aligned}[t]&\phantom{={}}19\times 40\\
&=40\times (20-1)\\
&=40\times 20-40\times 1\\
&=800-40\\
&=760\end{aligned}}

\MethodBlock{口算思路二：按数位拆分计算}{把一个因数拆成整十、整百数与余下部分之和。}{\begin{aligned}[t]&\phantom{={}}19\times 40\\
&=40\times (10+9)\\
&=40\times 10+40\times 9\\
&=400+360\\
&=760\end{aligned}}
\MethodTitle{第6小题：$360\div 20$}

\MethodBlock{口算思路一：同时去掉相同个数的0}{}{\begin{aligned}[t]&\phantom{={}}360\div 20\\
&=(360\div 10)\div (20\div 10)\\
&=36\div 2\\
&=18\end{aligned}}

\MethodBlock{口算思路二：拆分被除数}{}{\begin{aligned}[t]&\phantom{={}}360\div 20\\
&=(200+160)\div 20\\
&=200\div 20+160\div 20\\
&=10+8\\
&=18\end{aligned}}
\MethodTitle{第7小题：$630\div 90$}

\MethodBlock{口算思路一：同时去掉相同个数的0}{}{\begin{aligned}[t]&\phantom{={}}630\div 90\\
&=(630\div 10)\div (90\div 10)\\
&=63\div 9\\
&=7\end{aligned}}

\MethodBlock{口算思路二：拆分除数，连续相除}{}{\begin{aligned}[t]&\phantom{={}}630\div 90\\
&=630\div (2\times 45)\\
&=630\div 2\div 45\\
&=315\div 45\\
&=7\end{aligned}}
\MethodTitle{第8小题：$32\times 30$}

\MethodBlock{口算思路一：凑整后补偿}{把接近整十、整百的因数凑整，再补回或减去多算的部分。}{\begin{aligned}[t]&\phantom{={}}32\times 30\\
&=30\times (30+2)\\
&=30\times 30+30\times 2\\
&=900+60\\
&=960\end{aligned}}
\MethodTitle{第9小题：$20\times 15$}

\MethodBlock{口算思路一：按数位拆分计算}{把一个因数拆成整十、整百数与余下部分之和。}{\begin{aligned}[t]&\phantom{={}}20\times 15\\
&=20\times (10+5)\\
&=20\times 10+20\times 5\\
&=200+100\\
&=300\end{aligned}}
\MethodTitle{第10小题：$450\div 15$}

\MethodBlock{口算思路一：利用商不变的性质}{被除数和除数同时乘同一个非零数，商不变。}{\begin{aligned}[t]&\phantom{={}}450\div 15\\
&=(450\times 2)\div (15\times 2)\\
&=900\div 30\\
&=30\end{aligned}}

\MethodBlock{口算思路二：拆分除数，连续相除}{}{\begin{aligned}[t]&\phantom{={}}450\div 15\\
&=450\div (5\times 3)\\
&=450\div 5\div 3\\
&=90\div 3\\
&=30\end{aligned}}
\MethodTitle{第12小题：$500\div 25$}

\MethodBlock{口算思路一：利用商不变的性质}{被除数和除数同时乘同一个非零数，商不变。}{\begin{aligned}[t]&\phantom{={}}500\div 25\\
&=(500\times 2)\div (25\times 2)\\
&=1000\div 50\\
&=20\end{aligned}}

\MethodBlock{口算思路二：拆分除数，连续相除}{}{\begin{aligned}[t]&\phantom{={}}500\div 25\\
&=500\div (5\times 5)\\
&=500\div 5\div 5\\
&=100\div 5\\
&=20\end{aligned}}


\par\noindent

\QHeading{第2题　计算填表。}

\DataTable{4}{物品名称 & 单价/元 & 购买数量/件 & 总价/元\\\hline
衬衫 & \(5\) & \(36\) & {\color{answerred}180}\\\hline
裙子 & \(89\) & {\color{answerred}8} & \(712\)\\\hline
雨衣 & {\color{answerred}36} & \(18\) & \(648\)\\\hline
泳裤 & \(16\) & \(45\) & {\color{answerred}720}\\\hline}
\par

\par\noindent

\Answer{
\QHeading{计算过程：计算填表}
\begin{center}
\begin{adjustbox}{max width=\linewidth}\begin{tabular}{c|rrr}
\hline
物品名称&单价/元&购买数量/件&总价/元\\\hline
衬衫&5&36&180\\
裙子&89&8&712\\
雨衣&36&18&648\\
泳裤&16&45&720\\\hline
\end{tabular}\end{adjustbox}
\end{center}
\begin{enumerate}
\item 衬衫总价：\(\begin{aligned}[t]&\phantom{={}}5\times36\\&=5\times(30+6)\\&=5\times30+5\times6\\&=150+30\\&=180\text{（元）}\text{。}\end{aligned}\)
\item 裙子数量：\(712\div89=8\)件，核对\(89\times8=712\)。
\item 雨衣单价：\(648\div18=36\)元，核对\(36\times18=648\)。
\item 泳裤总价：\(\begin{aligned}[t]&\phantom{={}}16\times45\\&=16\times(40+5)\\&=16\times40+16\times5\\&=640+80\\&=720\text{（元）}\text{。}\end{aligned}\)
\end{enumerate}

}

\QHeading{第3题　解决问题。}

\Problem{①超市里一箱牛奶58元，一袋面包7元。妈妈要买3箱牛奶和1袋面包，一共需要花多少钱？}
\par


\MethodTitle{方法一}
\Answer{\Formula{\begin{aligned}[t]&\phantom{={}}58\times3+7\\&=174+7\\&=181\text{（元）}\end{aligned}}\par 答：一共需要花181元。}

\MethodTitle{方法二：按每箱60元估算后扣除多算的费用}
\Answer{每箱牛奶多算2元，3箱要减6元，再加面包钱。}
\Answer{\Formula{\begin{aligned}[t]&\phantom{={}}60\times 3-2\times 3+7\\
&=180-6+7\\
&=174+7\\
&=181\text{（元）}\end{aligned}}}
\Answer{答：一共需要花181元。}
\Problem{②文具店里每支钢笔12元，已知3支钢笔的总价和6支圆珠笔的总价一样，每支圆珠笔多少钱？}

\MethodTitle{方法一}
\Answer{\Formula{\begin{aligned}[t]&\phantom{={}}12\times3\div6\\&=36\div6\\&=6\text{（元）}\end{aligned}}\par 答：每支圆珠笔6元。}

\MethodTitle{方法二：先比较支数}
\Answer{3支钢笔的价钱等于6支圆珠笔的价钱，所以1支钢笔的价钱等于2支圆珠笔的价钱。}
\Answer{\Formula{\begin{aligned}[t]&\phantom{={}}12\div (6\div 3)\\
&=12\div 2\\
&=6\text{（元）}\end{aligned}}}
\Answer{答：每支圆珠笔6元。}

